% !PRINTED 217
\weberpage{217}
%<<CONT
equation for a smaller divisor. In contrast to this we call the
equation whose roots \gesperrt{are all} the $x_{\mu,\mu'}$ % [sic]
the \gesperrt{general division equation} for the divisor $n$.

By any substitution of the group $\mathfrak{A}$
\[
S = \begin{pmatrix} a, & b\\ c, & \partial \end{pmatrix}
\]
$(1,0)$, $(0,1)$ pass over into $(\partial, -c)$, $(-b, a)$, and so if $S$ is not
the identical substitution, then certainly at least one of the
two roots $x_{1,0}$, $x_{0,1}$ is altered by $S$. From this it results
\gesperrt{that the Galois group of the proper division equation
is exactly the same as that of the general one,}
only applied to the case in which $\mu, \mu', n$ have no common
divisor, namely, according as the $n$th roots of unity
are adjoined or not, $\mathfrak{B}$ or $\mathfrak{A}$.

From this it follows further \gesperrt{that the proper division equation
is irreducible, even after the adjunction of arbitrary constants.}

For if $\gamma, \delta$ are any two numbers having no common divisor
with $n$, then one can determine $\alpha, \beta$ in accordance with the congruence
\[
\alpha\delta - \beta\gamma \equiv 1 \pmod{n}
\]
and apply the substitution contained in $\mathfrak{B}$
\[
T = \begin{pmatrix} \alpha, & \beta\\ \gamma, & \delta \end{pmatrix}
\]
to every rational equation between the roots $x_{\mu,\mu'}$.
If then $x_{1,0}$ satisfies a rational equation (with arbitrary constant
coefficients)
\[
\Phi(x_{1,0}) = 0
\]
then every other root $x_{\delta,-\gamma}$ of the proper division equation
satisfies the same equation, from which the irreducibility of the
latter follows.

\begin{center}
\S\,65. \gesperrt{Reduction of the division equation}\\
\gesperrt{to transformation equations.}
\end{center}
The roots of the proper division equation can be arranged
in \gesperrt{rows} in the following way. Choose at will
one of the roots:
\[
x_{\mu_1,\mu'_1} = \mathrm{sn}\left(\frac{4\mu_1 K + 4\mu'_1 iK'}{n}\right) = \mathrm{sn}\,\Omega_1.
\]
